\section{Antecedentes}

El uso de videojuegos y simulaciones para enseñar ciberseguridad no es una propuesta aislada. Ng y Hasan (2024), en una revisión sistemática publicada en \textit{Computers \& Security}, analizaron el desarrollo de juegos serios de ciberseguridad y señalan que estos recursos se han aplicado en temas como conocimientos básicos de seguridad, redes, defensa, ataques éticos y entrenamiento de usuarios. Esta revisión es relevante porque muestra que los juegos serios no solo se han usado como recursos motivacionales, sino también como medios para organizar prácticas de aprendizaje.

Uno de los antecedentes clásicos es \textit{CyberCIEGE}, presentado por Irvine, Thompson y Allen (2005) en \textit{IEEE Security \& Privacy}. Este videojuego fue diseñado para enseñar conceptos de seguridad de la información mediante escenarios donde el usuario toma decisiones, administra recursos y observa consecuencias dentro de un entorno simulado. Su importancia para el presente proyecto está en demostrar que una dinámica de juego puede representar problemas de seguridad sin trabajar sobre sistemas reales.

También existen experiencias orientadas al entrenamiento práctico. Thompson y otros (2011) explican que \textit{CyberCIEGE} permite enfrentar al estudiante con puntos de decisión, experimentar fallos y reflexionar sobre sus resultados. Este enfoque se relaciona con el proyecto porque la intención no es que el estudiante memorice una lista de amenazas, sino que pueda practicar reconocimiento de señales, análisis de contexto y selección de respuestas.

En cuanto a la necesidad actual de formación en ciberseguridad, el marco NICE del NIST (2020) destaca la importancia de desarrollar conocimientos, habilidades y tareas asociadas al trabajo en ciberseguridad. Este enfoque refuerza la necesidad de que los estudiantes no solo conozcan amenazas, sino que practiquen acciones relacionadas con su identificación y respuesta. Por ello, una herramienta de simulación puede servir como apoyo para acercar al estudiante a situaciones similares a las del contexto profesional.

Finalmente, el informe DBIR de Verizon (2025) mantiene la relevancia del factor humano en las brechas de seguridad, incluyendo errores, ingeniería social y uso indebido de credenciales. Estos datos justifican que el proyecto incluya escenarios como \textit{phishing}, llamadas sospechosas y decisiones de respuesta ante incidentes, ya que son situaciones frecuentes en la realidad digital actual. A partir de estos antecedentes, el proyecto se orienta a construir una experiencia práctica, progresiva y segura para fortalecer habilidades de respuesta ante amenazas digitales.
